\chapter{Conclusión}

A lo largo de este trabajo, analizamos la herramienta \textit{Java Typestate Checker}, la cual es un verificador estático de protocolos de objetos basada en \textit{typestates}.

Comenzamos por mencionar trabajos relacionados y previos a la herramienta. Por un lado, la introducción inicial y el marco teórico de la abstracción \textit{typestate}: qué representan, cómo se componen, y cómo se pueden utilizar para modelar el estado de los objetos. Por otro lado, repasamos herramientas y lenguajes predecesores, las cuales ya colocaban a los \textit{typestate} en el núcleo de su dominio y ya contaban con funcionalidades de verificación de protocolos de objetos.

Luego, procedimos a dar una explicación alargada de \textit{JaTyC}. Abarcamos su próposito, su alcance y algunas de sus limitaciones conocidas. Expusimos sus características principales, entre las cuales se encontraban su sintaxis, formato de protocolos, anotaciones adicionales, manejo de aliasing y subtipado. También mencionamos cómo utiliza los \textit{typestates} para modelar sus protocolos.

Por otro lado, profundizamos en la implementación de su versión inicial. Expusimos cómo modela internamente a los diferentes actores involucrados en la verificación de protocolos, su sistema propio de tipos con reticulado, y su división operativa en dos fases: la de inferencia de tipos y la de verificación de errores. También escribimos sobre su relación con el framework sobre el cual está construida, el \textbf{\textit{Checker Framework}}. Mostramos qué clases y conceptos del mismo utilizaba, y su utilización de las entidades del compilador de Java.

Finalmente, mencionamos varios de los experimentos realizados con la herramienta. En un primer lugar, buscábamos ampliar nuestro conocimiento sobre sus capacidades, y su comportamiento en casos de uso y estructuras de código cotidianas. En un segundo lugar, intentamos descubrir sus limitaciones, situaciones en las que se comportase en forma distinta a lo esperado, y cómo afectaba la calidad y flexibilidad de los códigos producibles tales que no reporte errores. Por último, la probamos en casos bordes particulares, exponiendo debilidades o lugares de mejora a futuro.

Sobre los protocolos compuestos, investigamos cómo se comportaba \textit{JaTyC} cuando tanto una clase como uno de sus colaboradores internos tenían protocolos. Encontramos que realiza verificaciones modulares, donde el comportamiento del cliente del objeto se analizan por separado del comportamiento del objeto por sobre su colaborador interno. En ese sentido, \textit{JaTyC} no ve a los protocolos como compuestos, sino como dos entidades separadas las cuales no dependen entre si. A su vez, reconocimos su visión pesimista y potenciales falsos positivos, donde la herramienta reporta errores en programas los cuales pueden ser técnicamente válidos o correctos.

La siguiente sección relata sobre la diferencia entre la correctitud y el orden de los métodos. En ella notamos que la herramienta no tiene conocimiento sobre la especificación de correctitud de un método más allá de las anotaciones de \textit{typestates} con las que se los decora. Más bien, simplemente predica sobre el orden de llamado a los mismos. Esto implica que \textit{JaTyC} no asegura que un objeto o método esté implementado correctamente, y que esto queda en manos de otras herramientas de verificación de correctitud y el programador. A su vez, notamos que podemos mejorar la seguridad y correctitud de un programa a través de los protocolos de la herramienta, forzando a clientes a utilizar métodos de control antes de cualquier operación critica en objetos. En ese sentido, fomenta que los desarrolladores sean más defensivos al programar.

En el apartado de intervención humana, buscábamos entender qué mecanismos de control posee \textit{JaTyC} a la hora de verificar el diseño de protocolos, y cuanta influencia puede tener un programador en sus resultados. Luego de nuestros experimentos notamos que la herramienta es altamente susceptible al error humano, donde protocolos mal definidos pueden alterar gravemente el resultado del análisis. Cualquier protocolo puede ser ``sobreespecificado'' o ``subespecificado'', provocando que la herramienta reporte errores para cadenas de métodos válidas o ignore llamadas que querríamos que prohiba. En ese sentido, consideramos esta una gran debilidad de la herramienta: a cuanta más complejidad de un programa más puede elevarse la dificultad de correctamente definir los protocolos de las entidades, y por ende más crece la probabilidad de cometer errores que influyan en el resultado de la verificación. Mencionamos como posible mejora al sistema la generación automática de los protocolos, tomando como oráculo por ejemplo a al implementación actual del objeto y sus métodos de control.

En la sección de subtipado expusimos a la herramienta a distintas situaciones de subtipado, con diferentes configuraciones de clases con y sin protocolos. Encontramos que la misma requiere que subprotocolos incluyan a todas las cadenas de llamadas a métodos válidas de su protocolo padre, incluidas las transiciones a \texttt{end}. Los subprotocolos solo son capaces de agregar nuevas secuencias permitidas. Esto busca que subclases sigan comportandose como su superclase o interfaz, y que sus clientes tengan experiencias de uso consistentes.

Luego, se evaluó la diferencia entre el estado de un objeto, interpretado como los valores de sus colaboradores internos, y su \textit{typestate}. Buscamos entender cuánto tiene la herramienta en cuenta el estado de un objeto a la hora de realizar la verificación de protocolos, su posible influencia en los métodos disponibles de la instancia, y si existe alguna correlación entre un \textit{typestate} y un estado. Nuestra experimentación encontró que el análisis de \textit{JaTyC} se lleva a cabo sobre la abstracción de los \textit{typestates} en forma prácticamente independiente a los estados internos de los objetos. Verifica que métodos se llamen en correcto orden pero no que un \textit{typestate} particular se corresponda a una o varias configuraciones válidas de estados internos. Similar a lo denotado anteriormente, la herramienta no garantiza que los objetos se encuentren con estados internos válidos a la hora de llamar a sus métodos. Que cada \textit{typestate} se corresponda a estos estados internos correctos queda en manos de los desarrolladores.

En la siguiente sección se experimentó sobre las capacidades de modelado de la herramienta, donde buscamos ponerla a prueba en diferentes situaciones para encontrar sus límites y flexibilidad. Hallamos que da soporte a la sobrecarga de métodos y que permite el ``prestado'' de referencias lineales al llamar métodos. A su vez, notamos que no da soporte directo a suscribir clases a protocolos condicionalmente, util cuando se quiere especificar comportamientos distintos para objetos inmutables o constructores sobrecargados. De todos modos, esto puede en cierta forma solventarse modelando nuestros objetos y protocolos en formas particulares.

Revisamos también cómo la herramienta reacciona al aliasing en dos sitaciones particulares: cuando se delega la terminación de protocolo a un método anotado con \texttt{@Requires}, y cuando una clase tiene métodos \textit{getters} para colaboradores protocolados. Por un lado, encontramos que \textit{JaTyC} correctamente realiza el seguimiento de la finalización de protocolo al pasar referencias lineales como argumentos. Esto puede verse como una delegación de responsabilidades desde un método cliente hacia el método llamado. Por otro lado, vimos que entregar la referencia lineal de un colaborador interno provoca grandes cambios en los métodos permitidos de la clase y perjudica su reutilización. En ese sentido, \textit{JaTyC} parece fomentar otras formas de diseño de programas, donde estas referencias internas permanecen en los colaboradores internos de las clases, y las mismas exponen métodos útiles para que sus clientes puedan acceder a ellos.

Finalmente, notamos la limitación de \textit{JaTyC} a la hora de protocolar excepciones. En otras palabras, la herramienta permite transiciones entre \textit{typestates} a través de llamadas de métodos y sus resultados, pero no a través de excepciones esperadas. Esto flexibilizaría más el análisis de los protocolos de objetos, daría visibilidad a excepciones que deberían de suceder, y permitiría seguir manteniendo bajo control el análisis del estado de un objeto luego de que las excepciones esperadas se levanten.

En conclusión, consideramos que \textit{JaTyC} es una herramienta útil para especificar y verificar protocolos de objetos estáticamente en Java. A su vez, a lo largo de este trabajo identificamos diversas limitaciones que condicionan su alcance y flexibilidad. En particular, observamos que su análisis se encuentra fuertemente atado a la abstracción de los \textit{typestates}, sin considerar directamente la validez de los estados internos de los objetos ni la correctitud de las implementaciones de sus métodos. Asimismo, su dependencia de protocolos definidos manualmente introduce una fuente importante de errores, la cual podría volverse más relevante a mayor tamaño y complejidad de los programas. Finalmente, es incapaz de representar con sus protocolos excepciones esperadas y el estado de un objeto luego de levantarlas. Avances en estas areas podrían ampliar las capacidades y alcance de la herramienta.